\documentclass[a4paper,10pt]{article}
\usepackage[utf8]{inputenc}
\usepackage[T1]{fontenc}
\usepackage{geometry}
\usepackage{xcolor}
\usepackage{array}
\usepackage{longtable}
\usepackage{booktabs}
\usepackage{enumitem}
\usepackage{hyperref}
% Note: babel and microtype are not available on this build system

\geometry{margin=2.5cm}

\definecolor{high}{HTML}{E74C3C}
\definecolor{medium}{HTML}{F39C12}
\definecolor{low}{HTML}{27AE60}
\definecolor{info}{HTML}{2980B9}

\newcommand{\severity}[1]{%
  \ifstrequal{#1}{Critical}{\textcolor{high}{\textbf{#1}}}{%
  \ifstrequal{#1}{High}{\textcolor{high}{\textbf{#1}}}{%
  \ifstrequal{#1}{Medium}{\textcolor{medium}{\textbf{#1}}}{%
  \ifstrequal{#1}{Low}{\textcolor{low}{\textbf{#1}}}{#1}}}}}

\title{\textbf{Vulkan Engine Performance Evaluation Report}}
\author{Static Code Analysis}
\date{\today}

\begin{document}
\maketitle

\begin{abstract}
This report presents a comprehensive performance evaluation of the in-development Vulkan
graphics engine, based on static analysis of the source code. It identifies performance
bottlenecks, unsafe synchronization patterns, and API modernization opportunities. Each
critical issue is classified by severity and accompanied by mitigation recommendations.
\end{abstract}

\tableofcontents
\newpage

\section{Executive Summary}

The engine predominantly uses modern Vulkan with \textbf{dynamic rendering} across all
pipelines, which is a very positive aspect. However, \textbf{4 Critical/High severity
issues}, \textbf{5 Medium severity issues}, and \textbf{4 Low severity issues} were
identified. The most impactful problems involve per-frame \texttt{std::thread} creation,
heavy use of \texttt{vkCmdUpdateBuffer} (which causes implicit queue stalls), complete
absence of a pipeline cache, and not using VMA (Vulkan Memory Allocator) for images.

\vspace{0.5em}
\noindent
\textbf{Overall rating: \textcolor{medium}{Fair}} --- the engine works, but leaves
significant performance on the table due to poor CPU--GPU synchronization practices and
inefficient allocation patterns.

\section{Methodology}

The analysis was performed through:
\begin{itemize}[nosep]
    \item Complete source code review (C++, GLSL).
    \item Identification of Vulkan API calls (legacy vs.\ modern).
    \item Manual mapping of synchronization barriers, memory allocations, and resource
          creation.
    \item Search for known performance loss patterns (\textit{stalls}, \textit{bubbles},
          temporary allocations, mutex contention).
\end{itemize}

\section{Severity Classification}

\begin{itemize}[nosep]
    \item[\textbf{Critical}] --- Causes per-frame CPU--GPU stalls, severe FPS degradation,
          or corruption risk.
    \item[\textbf{High}] --- Significant impact on performance, scalability, or latency.
    \item[\textbf{Medium}] --- Moderate impact; incremental improvement.
    \item[\textbf{Low}] --- Marginal impact or only in rare situations.
\end{itemize}

\section{Identified Critical Issues}

\subsection{Per-Frame Thread Creation (Critical)}

\begin{tabular}{ll}
    \textbf{Severity} & \severity{Critical} \\
    \textbf{Files}     & \texttt{main.cpp:1124--1244} \\
    \textbf{Pattern}   & \texttt{std::thread} creation per frame in
                          \texttt{drawSidebar()} \\
\end{tabular}

\vspace{0.5em}
\noindent
\textbf{Description:} Whenever water is enabled, the code spawns a new thread
(\texttt{asyncTasks.emplace\_back(...)} at \texttt{main.cpp:1144}) for the back-face
pass. This thread is \textbf{created, executed, and destroyed every frame}. The
cost of creating a \texttt{std::thread} is hundreds of microseconds,
directly impacting the per-frame CPU budget. Additionally, duplicate joins
(\texttt{main.cpp:1249} and \texttt{1272}) indicate possible redundant waiting.

\noindent\textbf{Recommendation:} Replace with a fixed \textbf{thread pool} using
\texttt{std::jthread} or \texttt{std::async} with \texttt{launch::async} policy.
Keeping threads alive and reusing them reduces creation \textit{overhead} to
virtually zero.

\subsection{\texttt{vkCmdUpdateBuffer} --- Queue Stalls (Critical)}

\begin{tabular}{ll}
    \textbf{Severity} & \severity{Critical} \\
    \textbf{Files}     & \texttt{VegetationRenderer.cpp:422,427},
                          \texttt{SceneRenderer.cpp:236,323},
                          \texttt{IndirectRenderer.cpp:1030},
                          \texttt{Solid360Renderer.cpp:373,628} \\
    \textbf{Occurrences}& \textbf{11 calls} total \\
\end{tabular}

\vspace{0.5em}
\noindent
\textbf{Description:} \texttt{vkCmdUpdateBuffer} is a \textbf{legacy} function that
copies data from CPU to GPU inside a command buffer, but implicitly inserts a
\texttt{FULL\_QUEUE} barrier (\textit{top-of-pipe} to \textit{bottom-of-pipe}),
draining the entire graphics queue. This eliminates any parallelism between uploads
and shader execution. The most critical occurrences are:

\begin{itemize}[nosep]
    \item \texttt{VegetationRenderer.cpp:422,427} --- Upload of up to 256 draw
          commands per frame (workaround for a RADV bug).
    \item \texttt{SceneRenderer.cpp:236,323} --- Shadow and main UBO uploads,
          4 cascades each, per frame.
    \item \texttt{IndirectRenderer.cpp:1030} --- Visible counter reset before culling.
\end{itemize}

\noindent\textbf{Recommendation:} Replace with \textbf{persistently mapped host-visible
buffers} (using \texttt{VK\_MEMORY\_PROPERTY\_HOST\_COHERENT\_BIT}
or \texttt{VK\_MEMORY\_PROPERTY\_HOST\_CACHED\_BIT} with
\texttt{vkFlushMappedMemoryRanges}). Map once at initialization and update
data with \texttt{memcpy}. This eliminates stalls and reduces the number of
synchronization barriers.

\subsection{Missing Pipeline Cache (High)}

\begin{tabular}{ll}
    \textbf{Severity} & \severity{High} \\
    \textbf{Files}     & Entire codebase \\
    \textbf{Evidence} & No occurrence of \texttt{VkPipelineCache} or
                          \texttt{vkCreatePipelineCache} \\
\end{tabular}

\vspace{0.5em}
\noindent
\textbf{Description:} All calls to \texttt{vkCreateGraphicsPipelines} and
\texttt{vkCreateComputePipelines} use \texttt{VK\_NULL\_HANDLE} as the
pipeline cache. This means every application run recompiles shaders from scratch
(SPIR-V to GPU machine code). In engines with dozens of pipelines, this process
can take seconds. The code does not even attempt to load/save a disk cache.

\noindent\textbf{Recommendation:} Create a \texttt{VkPipelineCache} at
initialization, serialize it to disk on shutdown with
\texttt{vkGetPipelineCacheData}, and reload it on the next run.
Additionally, use \texttt{VK\_KHR\_pipeline\_binary} (Vulkan 1.4) if
available for more granular caching.

\subsection{VMA Not Used for Images (High)}

\begin{tabular}{ll}
    \textbf{Severity} & \severity{High} \\
    \textbf{Files}     & \texttt{VulkanApp.cpp:1602},
                          \texttt{RendererUtils.hpp:67},
                          \texttt{WaterRenderer.cpp:1079,1138},
                          \texttt{ImpostorCapture.cpp:450,510,571,634},
                          \texttt{Solid360Renderer.cpp:68},
                          \texttt{BillboardService.cpp:142},
                          \texttt{TextureArrayManager.cpp:271},
                          and 10+ more locations \\
    \textbf{Occurrences}& \textbf{25+ direct} \texttt{vkAllocateMemory} calls \\
\end{tabular}

\vspace{0.5em}
\noindent
\textbf{Description:} VMA (\textit{Vulkan Memory Allocator}) is used for
buffers (\texttt{VulkanApp.cpp:4305}, \texttt{createBuffer()}), but
\textbf{all images} are allocated via direct \texttt{vkAllocateMemory}.
This means each image has its own memory allocation, losing the benefits of
suballocation (fewer \texttt{vkAllocateMemory} calls, better memory utilization,
automatic defragmentation).

\noindent\textbf{Recommendation:} Extend VMA usage to images by creating
a \texttt{createImageWithVma()} function using
\texttt{vmaCreateImage()}. This reduces the number of memory allocations and
improves fragmentation. VMA is already a project dependency, so there is no
additional cost.

\subsection{Blocking Wait per Frame (High)}

\begin{tabular}{ll}
    \textbf{Severity} & \severity{High} \\
    \textbf{Files}     & \texttt{VulkanApp.cpp:4540} \\
    \textbf{Pattern}   & Blocking \texttt{vkWaitForFences} in \texttt{drawFrame()} \\
\end{tabular}

\vspace{0.5em}
\noindent
\textbf{Description:} At the start of each frame, the CPU blocks with
\texttt{vkWaitForFences(device, 1, \&inFlightFences[currentFrame], VK\_TRUE,
UINT64\_MAX)}. If the GPU is still processing the previous frame (or the oldest
in-flight frame), the CPU stays \textbf{idle} until the GPU finishes. With 3
frames in flight, this blocking is short for most frames, but when there is a
queue \textit{bubble} (e.g., heavy shadows), the CPU waits.

\noindent\textbf{Recommendation:} Maintaining \textbf{3 frames in flight} already
helps. To avoid stalls, consider introducing \textbf{timeline semaphores}
(\texttt{VK\_SEMAPHORE\_TYPE\_TIMELINE}) with \texttt{vkWaitSemaphores} to
wait only for the necessary execution point instead of the entire queue. This
enables finer CPU--GPU overlap.

\subsection{Temporary Command Pool per Async Submission (Medium)}

\begin{tabular}{ll}
    \textbf{Severity} & \severity{Medium} \\
    \textbf{Files}     & \texttt{VulkanApp.cpp:2512} \\
    \textbf{Pattern}   & \texttt{vkCreateCommandPool} inside
                          \texttt{submitCommandBufferAsyncToQueue()} \\
\end{tabular}

\vspace{0.5em}
\noindent
\textbf{Description:} Every async submission (including the back-face pass
in \texttt{main.cpp}) creates and destroys a temporary command pool.
Creating a command pool is not extremely expensive, but with dozens of async
submissions per second (water, vegetation, etc.), the \textit{overhead}
accumulates.

\noindent\textbf{Recommendation:} Pre-allocate a set of command pools
(\textit{pool ring}) at startup, one per submission thread, and reset them
with \texttt{vkResetCommandPool} instead of creating/destroying. This reduces
pressure on the Vulkan \textit{driver}.

\subsection{Redundant Chain Barriers (Medium)}

\begin{tabular}{ll}
    \textbf{Severity} & \severity{Medium} \\
    \textbf{Files}     & \texttt{ShadowRenderer.cpp},
                          \texttt{SceneRenderer.cpp} \\
    \textbf{Pattern}   & Multiple layout transitions per frame
\end{tabular}

\vspace{0.5em}
\noindent
\textbf{Description:} \texttt{ShadowRenderer.cpp} executes \textbf{9 layout
transitions} per frame (for 4 shadow cascades + blur + restore). \texttt{SceneRenderer.cpp}
adds \textbf{10 barriers} for UBOs (2 per cascade + 2 for the main UBO). Many
of these barriers could be combined or eliminated. Each barrier has a
non-negligible driver processing cost.

\noindent\textbf{Recommendation:} Group layout transitions using
\texttt{VkDependencyInfo} with multiple barriers in a single
\texttt{vkCmdPipelineBarrier2} call. For UBOs, the transition to
\texttt{TRANSFER\_WRITE} and back to \texttt{UNIFORM\_READ} can be
eliminated if UBOs are host-visible mapped (see Section~4.2).

\subsection{Shaders Compiled for Vulkan 1.1 (Medium)}

\begin{tabular}{ll}
    \textbf{Severity} & \severity{Medium} \\
    \textbf{Files}     & \texttt{Makefile:144} \\
    \textbf{Pattern}   & \texttt{glslc --target-env=vulkan1.1} \\
\end{tabular}

\vspace{0.5em}
\noindent
\textbf{Description:} The \texttt{Makefile} compiles all shaders with
\texttt{--target-env=vulkan1.1}. This means shaders are generated as
SPIR-V 1.3 (maximum supported by Vulkan 1.1), missing SPIR-V 1.6 capabilities
like \texttt{ScalarBlockLayout} (better uniform packing),
\texttt{MaximallyReconverges} (warp divergence optimization), and the
Vulkan 1.3+ memory model.

\noindent\textbf{Recommendation:} Update the target to
\texttt{--target-env=vulkan1.4} (or \texttt{vulkan1.3} as fallback). This
allows the GLSL compiler to generate more optimized SPIR-V. It is a one-line
change in the \texttt{Makefile}.

\subsection{Legacy \texttt{vkQueueSubmit} Usage (Medium)}

\begin{tabular}{ll}
    \textbf{Severity} & \severity{Medium} \\
    \textbf{Files}     & \texttt{VulkanApp.cpp:1420,1517,2131,2265,2348,4843} \\
    \textbf{Pattern}   & \texttt{vkQueueSubmit} (legacy) instead of
                          \texttt{vkQueueSubmit2} (modern) \\
\end{tabular}

\vspace{0.5em}
\noindent
\textbf{Description:} All queue submissions use \texttt{vkQueueSubmit} from the
original Vulkan API. The \texttt{vkQueueSubmit2} API (Vulkan 1.3+
\texttt{synchronization2}) offers a more expressive and efficient
synchronization model, allowing specification of \texttt{VkSemaphoreSubmitInfo}
and \texttt{VkCommandBufferSubmitInfo} with clearer semantics. Although the
practical difference is small, migration aligns the code with modern practices
and may benefit newer \textit{drivers}.

\noindent\textbf{Recommendation:} Gradually migrate to
\texttt{vkQueueSubmit2}. This requires restructuring the submission
structures to use \texttt{VkSubmitInfo2}, \texttt{VkSemaphoreSubmitInfo}, and
\texttt{VkCommandBufferSubmitInfo}. The change is mechanical, does not alter
semantics, and paves the way for timeline semaphores.

\subsection{Per-Node Octree Threads (Low)}

\begin{tabular}{ll}
    \textbf{Severity} & \severity{Low} \\
    \textbf{Files}     & \texttt{IteratorHandler.cpp:161--186} \\
    \textbf{Pattern}   & \texttt{std::thread} creation per child node
\end{tabular}

\vspace{0.5em}
\noindent
\textbf{Description:} In octree processing, up to 8 threads are created per node
to process children in parallel. For an octree with depth 8, this can result in
thousands of threads created/destroyed. Since this is not a per-frame operation
(it is used during construction/tessellation), the impact is reduced, but the
pattern is wasteful.

\noindent\textbf{Recommendation:} Replace with a global \textbf{thread pool}
(which also serves point 4.1) and distribute work with
\texttt{std::atomic} or concurrent task queues.

\subsection{\texttt{vkDeviceWaitIdle} in Initialization Paths (Low)}

\begin{tabular}{ll}
    \textbf{Severity} & \severity{Low} \\
    \textbf{Files}     & \texttt{VegetationRenderer.cpp:1093},
                          \texttt{VulkanResourceManager.cpp:331,353} \\
    \textbf{Pattern}   & \texttt{vkDeviceWaitIdle} in one-shot operations
\end{tabular}

\vspace{0.5em}
\noindent
\textbf{Description:} \texttt{vkDeviceWaitIdle} is called during \texttt{setImpostorData()}
(\texttt{VegetationRenderer.cpp:1093}) and during \texttt{resource manager}
destruction (\texttt{VulkanResourceManager.cpp:331,353}). Although these are
not in the per-frame hot path, they drain all queues and prevent any
parallelism.

\noindent\textbf{Recommendation:} Whenever possible, replace with targeted
waiting using \texttt{vkWaitForFences} on the relevant submission set. For
destruction, a full drain may be acceptable, but for impostor initialization,
a specific fence is more appropriate.

\subsection{Serial Shader Compilation (Low)}

\begin{tabular}{ll}
    \textbf{Severity} & \severity{Low} \\
    \textbf{Files}     & \texttt{Makefile:131--150} \\
    \textbf{Pattern}   & Serial shader compilation
\end{tabular}

\vspace{0.5em}
\noindent
\textbf{Description:} Shaders are compiled one at a time via Make's default
pattern rules. With 40+ shaders and dependencies between them (includes),
compilation time could be reduced.

\noindent\textbf{Recommendation:} Using \texttt{make -j} already applies
parallelism, but the Makefile itself could specify explicit dependencies to
maximize it. This is a minor improvement but a welcome one.

\section{Summary and Priorities}

\begin{longtable}{p{0.35\textwidth}ccp{0.25\textwidth}}
    \toprule
    \textbf{Critical Point} & \textbf{Severity} & \textbf{Effort} & \textbf{Impact} \\
    \midrule
    \endhead
    \bottomrule
    \endfoot
    Threads per frame
        & \severity{Critical}
        & Medium
        & Eliminates per-frame thread create/destroy \\
    \texttt{vkCmdUpdateBuffer}
        & \severity{Critical}
        & Low
        & Replace with mapped buffers; stalls eliminated \\
    Pipeline cache
        & \severity{High}
        & Low
        & Shaders compiled once; faster startup \\
    VMA for images
        & \severity{High}
        & Medium
        & Fewer allocations; better fragmentation \\
    Blocking wait
        & \severity{High}
        & Medium
        & Finer CPU--GPU overlap \\
    Temporary command pool
        & \severity{Medium}
        & Low
        & Reduced driver overhead \\
    Redundant barriers
        & \severity{Medium}
        & Medium
        & Less barrier processing \\
    Shaders Vulkan 1.1
        & \severity{Medium}
        & Minimal
        & Better shader optimization \\
    Legacy \texttt{vkQueueSubmit}
        & \severity{Medium}
        & High
        & Modernization; enables timeline semaphores \\
    Octree threads
        & \severity{Low}
        & Medium
        & Improvement in offline operations \\
    \texttt{vkDeviceWaitIdle}
        & \severity{Low}
        & Low
        & Avoids unnecessary full drains \\
    Serial compilation
        & \severity{Low}
        & Minimal
        & Marginally faster builds \\
\end{longtable}

\section{Immediate Recommendations}

Based on the \textbf{effort/impact} ratio, the following implementation order is
recommended:

\begin{enumerate}
    \item \textbf{Replace \texttt{vkCmdUpdateBuffer} with mapped buffers}
          (Section~4.2) --- low effort, eliminates coarse stalls.
    \item \textbf{Add a pipeline cache} (Section~4.3) --- low effort,
          immediate benefit in startup time.
    \item \textbf{Introduce a thread pool} (Section~4.1) --- medium effort,
          eliminates per-frame thread create/destroy.
    \item \textbf{Extend VMA to images} (Section~4.4) --- medium effort,
          positive impact on memory management.
    \item \textbf{Update shader target to Vulkan 1.4}
          (Section~4.8) --- minimal effort, trivial modernization.
\end{enumerate}

\section{Conclusion}

The engine demonstrates a solid foundation with correct use of
\textit{dynamic rendering} and modern synchronization patterns (\texttt{VkDependencyInfo},
\texttt{vkCmdPipelineBarrier2}). However, four categories of problems
significantly impact performance:

\begin{enumerate}[nosep]
    \item \textbf{CPU--GPU bottlenecks}: \texttt{vkCmdUpdateBuffer} and per-frame
          thread creation.
    \item \textbf{Inefficient allocation}: VMA not used for images, missing
          pipeline cache.
    \item \textbf{Excessive synchronization}: Blocking waits and redundant
          barriers.
    \item \textbf{Incomplete modernization}: Legacy \texttt{vkQueueSubmit} and
          shaders compiled for Vulkan 1.1.
\end{enumerate}

Fixing the Critical and High severity issues should result in a measurable FPS
improvement, especially in scenes with active water and vegetation. The engine's
overall rating is \textbf{Fair}, with the potential to evolve into \textbf{Good}
by implementing the above recommendations.

\end{document}
